\section{Avance 3: Juego Serio 3D en Godot 4.7.2, Fallas Diegéticas y Tablet SCADA In-Game}\label{sec:diseno_godot}

El Avance~3 consolida los modelos electromecánicos y térmicos del Avance~1 junto con los clasificadores inteligentes y la ruta de Mantenimiento Centrado en Confiabilidad (RCM) del Avance~2 dentro de un \emph{Juego Serio} (\textit{Serious Game}) 3D inmersivo desarrollado en \textbf{Godot Engine~4.7.2} sobre la tubería gráfica \textbf{Vulkan~1.3 Forward+}~\cite{GuiaProyecto2026,Singh2023}. Bajo el paradigma de la Industria~5.0~\cite{Breque2021,Xu2021}, la plataforma sitúa al estudiante como operario dentro del complejo agroindustrial de AgroMolinos del Valle S.A., integrando operación en primera persona, inducción normativa RETIE/LOTO, cinco perturbaciones físicas diegéticas sin botones artificiales y una Tablet Industrial SCADA completamente nativa en el motor.

\subsection{Arquitectura del Simulador 3D, Activos y Entorno Industrial}\label{subsec:arquitectura_godot}

% Figura 1: Arquitectura integral del proyecto (Transición Industria 4.0 -> Industria 5.0)
% AgroMolinos del Valle S.A. — Diseño por construcción (Layout by Construction)
% Requiere en preámbulo: IEEEtran, babel[spanish], amsmath, amssymb, siunitx,
% tikz con librerías: arrows.meta, positioning, calc, shapes.geometric, fit, backgrounds.
\begin{figure*}[!t]
\centering
\resizebox{0.96\textwidth}{!}{%
\begin{tikzpicture}[
    % Paleta sobria estilo IEEE (definida localmente para evitar colisiones)
    /utils/exec={%
        % Capa Base — Ámbar / Ingeniería Offline (Industria 4.0)
        \definecolor{archAmberDark}{HTML}{9A3412}%
        \definecolor{archAmberMed}{HTML}{C2410C}%
        \definecolor{archAmberBg}{HTML}{FFF7ED}%
        % Capa Central — Azul Cobalto / Hero (Godot 4.7.2)
        \definecolor{archBlueDark}{HTML}{1E3A8A}%
        \definecolor{archBlueMed}{HTML}{1D4ED8}%
        \definecolor{archBlueBg}{HTML}{EFF6FF}%
        % Bus Central & Edge AI — Violeta Técnico
        \definecolor{archVioletDark}{HTML}{5B21B6}%
        \definecolor{archVioletMed}{HTML}{6D28D9}%
        \definecolor{archVioletBg}{HTML}{F5F3FF}%
        % Mundo 3D — Verde Azulado Industrial
        \definecolor{archTealDark}{HTML}{0F766E}%
        % Capa Superior — Verde Esmeralda (Industria 5.0)
        \definecolor{archGreenDark}{HTML}{065F46}%
        \definecolor{archGreenMed}{HTML}{047857}%
        \definecolor{archGreenBg}{HTML}{ECFDF5}%
        % Neutros tipográficos
        \definecolor{archSlateDark}{HTML}{0F172A}%
        \definecolor{archSlateText}{HTML}{334155}%
        \definecolor{archSlateLine}{HTML}{CBD5E1}%
        \definecolor{archPanelBg}{HTML}{F8FAFC}%
    },
    font=\sffamily,
    % Estilos base de tarjetas: encabezado fijo en top mediante path picture + parbox[t]
    card base/.style 2 args={
        rectangle,
        rounded corners=3.5pt,
        draw=#1,
        line width=0.85pt,
        fill=white,
        inner xsep=5.5pt,
        inner ysep=4pt,
        align=left,
        text=archSlateDark,
        font=\sffamily\footnotesize,
        path picture={
            \fill[#2] (path picture bounding box.north west)
                rectangle ([yshift=-0.47cm]path picture bounding box.north east);
        }
    },
    card l1/.style={card base={archAmberMed}{archAmberDark}},
    card l2 phys/.style={card base={archBlueMed}{archBlueDark}},
    card l2 world/.style={card base={archTealDark}{archTealDark}},
    card l2 ai/.style={card base={archVioletMed}{archVioletDark}},
    card l3/.style={card base={archGreenMed}{archGreenDark}},
    % Contenedores de capa (en capa background)
    zone l1/.style={
        rectangle, rounded corners=5pt,
        draw=archAmberMed, line width=0.95pt,
        fill=archAmberBg,
        inner xsep=0.32cm, inner ysep=0.28cm
    },
    zone l2/.style={
        rectangle, rounded corners=5.5pt,
        draw=archBlueDark, line width=1.45pt,
        fill=archBlueBg,
        inner xsep=0.32cm, inner ysep=0.28cm
    },
    zone l3/.style={
        rectangle, rounded corners=5pt,
        draw=archGreenMed, line width=0.95pt,
        fill=archGreenBg,
        inner xsep=0.32cm, inner ysep=0.28cm
    },
    % Bus central dentro de Capa 2
    bus container/.style={
        rectangle, rounded corners=4pt,
        draw=archVioletDark, line width=1.05pt,
        fill=archVioletBg,
        inner xsep=0pt, inner ysep=6pt
    },
    bus pill/.style={
        rectangle, rounded corners=3pt,
        draw=archVioletMed, line width=0.8pt,
        fill=white,
        minimum height=0.64cm,
        inner xsep=7pt, inner ysep=2.5pt,
        align=center,
        font=\sffamily\footnotesize,
        text=archSlateDark
    },
    % Columna izquierda de progresión I4.0 -> I5.0
    spine card/.style 2 args={
        rectangle, rounded corners=4pt,
        draw=#1, line width=1.05pt,
        fill=#2,
        text width=2.72cm,
        inner xsep=5pt, inner ysep=6pt,
        align=center,
        font=\sffamily\footnotesize,
        text=white
    },
    % Estilos canónicos de flechas
    flow amber/.style={
        -{Stealth[length=5.5pt, width=4.2pt]},
        draw=archAmberDark, line width=1.1pt,
        shorten >=1.5pt, shorten <=1pt
    },
    flow blue bi/.style={
        {Stealth[length=4.8pt, width=3.5pt]}-{Stealth[length=4.8pt, width=3.5pt]},
        draw=archBlueDark, line width=1.0pt,
        shorten >=1.2pt, shorten <=1.2pt
    },
    flow green bi/.style={
        {Stealth[length=5.2pt, width=3.8pt]}-{Stealth[length=5.2pt, width=3.8pt]},
        draw=archGreenDark, line width=1.05pt,
        shorten >=1.5pt, shorten <=1.5pt
    },
    flow green uni/.style={
        -{Stealth[length=5.5pt, width=4.0pt]},
        draw=archGreenDark, line width=1.05pt,
        shorten >=1.5pt, shorten <=1pt
    },
    flow spine/.style={
        -{Stealth[length=6.5pt, width=5.2pt]},
        draw=archSlateDark, line width=1.5pt,
        shorten >=2pt, shorten <=2pt
    },
    edge lbl/.style={
        font=\sffamily\scriptsize,
        text=archSlateText,
        align=left,
        inner sep=1.5pt
    }
]
    % Parámetros maestros de espaciado (única fuente de verdad)
    \def\colw{5.15cm}
    \def\colgap{0.46cm}
    \def\busvgap{0.84cm}
    \def\layervgap{1.36cm}
    \def\spinegap{0.46cm}
    % Alturas internas exactas por fila (ancladas arriba vía \parbox[t][H][t])
    \def\hcardLtwo{2.42cm}
    \def\hcardLthree{2.34cm}
    \def\hcardLone{2.08cm}
    % Macro local para evitar partición de palabras y fijar altura idéntica en cada fila
    \def\archcard#1#2#3#4{%
        \parbox[t][#1][t]{\linewidth}{%
            \raggedright\hyphenpenalty=10000\exhyphenpenalty=10000%
            \strut\textbf{\color{white}#2}\\[5pt]%
            {\scriptsize\color{#3}#4}%
        }%
    }

    % =========================================================================
    % 1. CAPA CENTRAL (HERO) — NÚCLEO SERIOUS GAME GODOT 4.7.2 (VULKAN 1.3)
    % =========================================================================
    \node[card l2 phys, text width=\colw] (eng_phys) at (0,0) {%
        \archcard{\hcardLtwo}{Gemelo Físico}{archBlueDark}{%
            \texttt{InductionMotorKloss.gd}\\
            \texttt{VFDDriveController.gd}\\[3.5pt]
            \color{archSlateDark}\footnotesize
            $\bullet$ Modelo de Kloss (10\,HP, 460\,V).\\
            $\bullet$ Control escalar V/f + lazo PID.\\
            $\bullet$ $\text{kWh/ton}$ y $\text{CO}_2\text{e}$ (VFD vs. DOL).%
        }%
    };

    \node[card l2 world, text width=\colw, right=\colgap of eng_phys] (eng_3d) {%
        \archcard{\hcardLtwo}{Mundo 3D y Logística}{archTealDark}{%
            \texttt{planta\_completa.glb}\; (1.34M tris)\\[4.5pt]
            \color{archSlateDark}\footnotesize
            $\bullet$ 846 \texttt{BoxColliders} y 28 puertas PCA.\\
            $\bullet$ 15 máquinas industriales animadas.\\
            $\bullet$ Cadena de granos y logística de\\
            \phantom{$\bullet$ }camiones articulados en planta.%
        }%
    };

    \node[card l2 ai, text width=\colw, right=\colgap of eng_3d] (eng_ai) {%
        \archcard{\hcardLtwo}{Edge AI Nativo}{archVioletDark}{%
            \texttt{EdgeAIEngine.gd}\; (inferencia local)\\[4.5pt]
            \color{archSlateDark}\footnotesize
            $\bullet$ Red MLP (\textbf{316 pesos}) + SVM RBF.\\
            $\bullet$ Salud del motor: índice MHI y RUL.\\
            $\bullet$ Clasificación de fallas en vivo.\\
            $\bullet$ Motor de prescripción y alertas RCM.%
        }%
    };

    % Bus Central abarcando exactamente de eng_phys.west a eng_ai.east (misma línea Y)
    \coordinate (bus_left_anchor)  at ($(eng_phys.north west) + (0, \busvgap + 0.38cm)$);
    \coordinate (bus_right_anchor) at ($(eng_ai.north east)   + (0, \busvgap + 0.38cm)$);

    \node[
        rectangle, rounded corners=2.5pt,
        fill=archVioletDark, text=white,
        minimum height=0.64cm,
        font=\sffamily\scriptsize\bfseries, inner xsep=7pt, inner ysep=2.5pt,
        anchor=west
    ] (bus_tag) at ($(bus_left_anchor) + (0.18cm, 0)$) {\strut BUS CENTRAL};

    \node[bus pill, anchor=west] (bus_eb) at ($(bus_tag.east) + (0.45cm, 0)$) {%
        \strut\textbf{\texttt{EventBus.gd}} \;{\scriptsize\color{archVioletDark}(16 señales desacopladas)}%
    };

    \node[bus pill, anchor=east] (bus_sm) at ($(bus_right_anchor) + (-0.18cm, 0)$) {%
        \strut\textbf{\texttt{SimulationManager.gd}} \;{\scriptsize\color{archVioletDark}(bucle 10\,Hz)}%
    };

    % Contenedor invisible para calcular el bounding box exacto del Bus Central
    \node[inner xsep=0pt, inner ysep=6pt, fit=(bus_left_anchor)(bus_right_anchor)(bus_tag)(bus_eb)(bus_sm)] (bus_bar) {};

    % Contenedor de Capa Central primero en background, luego fondo del Bus Central encima
    \begin{pgfonlayer}{background}
        \node[zone l2, fit=(eng_phys)(eng_3d)(eng_ai)(bus_bar)] (zone_l2) {};
        \node[bus container, inner sep=0pt, fit=(bus_bar)] (bus_bar_bg) {};
    \end{pgfonlayer}

    % Flecha interna estrictamente horizontal del Bus Central
    \draw[{Stealth[length=4.5pt, width=3.4pt]}-{Stealth[length=4.5pt, width=3.4pt]},
          draw=archVioletDark, line width=1.0pt, shorten >=2pt, shorten <=2pt]
        (bus_eb.east) -- node[above=1pt, font=\sffamily\scriptsize\bfseries, text=archVioletDark] {Sync 10\,Hz} (bus_sm.west);

    % Conexiones verticales internas (Motores <-> Bus Central), todas idéntica longitud
    \draw[flow blue bi] (eng_phys.north) -- node[edge lbl, right=4pt] {Estado $I, T, \omega$} (bus_bar.south -| eng_phys.north);
    \draw[flow blue bi] (eng_3d.north)   -- node[edge lbl, right=4pt] {Escena / Colis.}    (bus_bar.south -| eng_3d.north);
    \draw[flow blue bi] (eng_ai.north)   -- node[edge lbl, right=4pt] {MHI / RUL / RCM}    (bus_bar.south -| eng_ai.north);

    % =========================================================================
    % 2. CAPA SUPERIOR — OPERACIÓN CENTRADA EN EL HUMANO (INDUSTRIA 5.0)
    % =========================================================================
    \node[card l3, text width=\colw, anchor=south]
        (hum_dieg) at ($(eng_phys.north |- zone_l2.north) + (0, \layervgap)$) {%
        \archcard{\hcardLthree}{Interacción Diegética 3D}{archGreenDark}{%
            \textit{Acción física sin botones 2D de falla}\\[4.5pt]
            \color{archSlateDark}\footnotesize
            $\bullet$ 5 fallas físicas inducibles en planta.\\
            $\bullet$ Vestuario EPP completo (6/6 RETIE).\\
            $\bullet$ 5 Reglas de Oro LOTO en campo.\\
            $\bullet$ Recorrido guiado en 20 estaciones.%
        }%
    };

    \node[card l3, text width=\colw, anchor=south]
        (hum_tab) at ($(eng_3d.north |- zone_l2.north) + (0, \layervgap)$) {%
        \archcard{\hcardLthree}{Tablet SCADA In-Game \texttt{[TAB]}}{archGreenDark}{%
            \textit{Nativa en Godot (sin navegador externo)}\\[4.5pt]
            \color{archSlateDark}\footnotesize
            $\bullet$ 9 pestañas: IoT, Cámaras 3D, RETIE\\
            \phantom{$\bullet$ }e IA; Rúbrica (0--5) e Historial.\\
            $\bullet$ Alertas push y diagnóstico en vivo.\\
            $\bullet$ Guías interactivas de mantenimiento.%
        }%
    };

    \node[card l3, text width=\colw, anchor=south]
        (hum_qa) at ($(eng_ai.north |- zone_l2.north) + (0, \layervgap)$) {%
        \archcard{\hcardLthree}{Certificación QA}{archGreenDark}{%
            \texttt{headless\_verify.gd}\; (CI determinística)\\[4.5pt]
            \color{archSlateDark}\footnotesize
            $\bullet$ Suite de pruebas: \textbf{10/10 PASS}.\\
            $\bullet$ Verificación de física e inferencia IA.\\
            $\bullet$ Auditoría de flujo de seguridad.\\
            $\bullet$ Exportación de reporte \texttt{JSON} RETIE.%
        }%
    };

    \begin{pgfonlayer}{background}
        \node[zone l3, fit=(hum_dieg)(hum_tab)(hum_qa)] (zone_l3) {};
    \end{pgfonlayer}

    % Flechas verticales Capa Central (Bus Central) <-> Capa Superior
    \draw[flow green bi] (bus_bar.north -| hum_dieg.south) -- (hum_dieg.south);
    \node[edge lbl, anchor=west] at ($(zone_l2.north -| hum_dieg.south)!0.5!(zone_l3.south -| hum_dieg.south) + (0.12cm, 0)$)
        {5 fallas / EPP / LOTO};

    \draw[flow green bi] (bus_bar.north -| hum_tab.south) -- (hum_tab.south);
    \node[edge lbl, anchor=west] at ($(zone_l2.north -| hum_tab.south)!0.5!(zone_l3.south -| hum_tab.south) + (0.12cm, 0)$)
        {Telemetría 10\,Hz / Push};

    \draw[flow green uni] (bus_bar.north -| hum_qa.south) -- (hum_qa.south);
    \node[edge lbl, anchor=west] at ($(zone_l2.north -| hum_qa.south)!0.5!(zone_l3.south -| hum_qa.south) + (0.12cm, 0)$)
        {Trazas / Log JSON};

    % =========================================================================
    % 3. CAPA BASE — AVANCES 1 Y 2 (MATLAB/SIMULINK & ENTRENAMIENTO IA)
    % =========================================================================
    \node[card l1, text width=7.26cm, anchor=north west]
        (a1) at ($(eng_phys.south west |- zone_l2.south) + (0, -\layervgap)$) {%
        \archcard{\hcardLone}{Avance 1: MATLAB / Simulink}{archAmberDark}{%
            \texttt{InductionMachineDTC.slx}\; $\cdot$\; Motor 10\,HP / 460\,V\\[4.5pt]
            \color{archSlateDark}\footnotesize
            $\bullet$ Modelado electromecánico: lazo PI + control DTC.\\
            $\bullet$ Red térmica equivalente de Cauer (estator y rotor).\\
            $\bullet$ Severidad de vibración normalizada ISO~10816.%
        }%
    };

    \node[card l1, text width=7.26cm, anchor=north east]
        (a2) at ($(eng_ai.south east |- zone_l2.south) + (0, -\layervgap)$) {%
        \archcard{\hcardLone}{Avance 2: Entrenamiento IA y RCM}{archAmberDark}{%
            Dataset 1200 reg. ($\mathbf{x}\in\mathbb{R}^8$)\; $\cdot$\; App Designer y Ruta RCM\\[4.5pt]
            \color{archSlateDark}\footnotesize
            $\bullet$ MLP $8\to 16\to 8\to 4$ + SVM RBF (\textbf{100\,\% test}).\\
            $\bullet$ Vector $\mathbf{x}\in\mathbb{R}^8$: rasgos eléctricos, térmicos y vibratorios.\\
            $\bullet$ Matriz de criticidad y plan de mantenimiento RCM.%
        }%
    };

    \begin{pgfonlayer}{background}
        \node[zone l1, fit=(a1)(a2)] (zone_l1) {};
    \end{pgfonlayer}

    % Flecha horizontal A1 -> A2 (Dataset)
    \draw[flow amber] (a1.east) -- (a2.west);
    \node[edge lbl, align=center, anchor=south] at ($(a1.east)!0.5!(a2.west) + (0, 0.10cm)$)
        {Dataset\\1200 reg.\\$\mathbf{x}\in\mathbb{R}^8$};

    % Flechas verticales Capa Base -> Capa Central (estrictamente verticales)
    \draw[flow amber] (a1.north -| eng_phys.south) -- (eng_phys.south);
    \node[edge lbl, anchor=west] at ($(zone_l1.north -| eng_phys.south)!0.5!(zone_l2.south -| eng_phys.south) + (0.12cm, 0)$)
        {Parámetros 10\,HP / Kloss / Cauer};

    \draw[flow amber] (a2.north -| eng_ai.south) -- (eng_ai.south);
    \node[edge lbl, anchor=east] at ($(zone_l1.north -| eng_ai.south)!0.5!(zone_l2.south -| eng_ai.south) + (-0.12cm, 0)$)
        {316 pesos MLP + SVM + RCM};

    % =========================================================================
    % 4. COLUMNA LATERAL DE PROGRESIÓN (INDUSTRIA 4.0 -> INDUSTRIA 5.0)
    % =========================================================================
    \node[spine card={archAmberDark}{archAmberDark}, minimum height=2.22cm, anchor=east]
        (sp_l1) at ($(zone_l1.west) + (-\spinegap, 0)$) {%
        {\scriptsize\textbf{CAPA BASE}}\\[3pt]
        {\small\textbf{Industria 4.0}}\\[3pt]
        {\scriptsize Avances 1 y 2\\MATLAB/Simulink\\y Entrenamiento IA}%
    };

    \node[spine card={archBlueDark}{archBlueDark}, minimum height=3.05cm, anchor=east]
        (sp_l2) at ($(zone_l2.west) + (-\spinegap, 0)$) {%
        {\scriptsize\textbf{CAPA CENTRAL}}\\[3pt]
        {\small\textbf{Godot 4.7.2}}\\[3pt]
        {\scriptsize Vulkan 1.3 Forward+\\Núcleo Serious Game\\Tiempo Real (10\,Hz)}%
    };

    \node[spine card={archGreenDark}{archGreenDark}, minimum height=2.36cm, anchor=east]
        (sp_l3) at ($(zone_l3.west) + (-\spinegap, 0)$) {%
        {\scriptsize\textbf{CAPA SUPERIOR}}\\[3pt]
        {\small\textbf{Industria 5.0}}\\[3pt]
        {\scriptsize Operación Centrada\\en el Humano y\\Certificación QA}%
    };

    % Conectores laterales limpios entre columna de progresión y zonas
    \draw[draw=archAmberMed, line width=1.1pt] (sp_l1.east) -- (zone_l1.west);
    \draw[draw=archBlueDark, line width=1.3pt] (sp_l2.east) -- (zone_l2.west);
    \draw[draw=archGreenMed, line width=1.1pt] (sp_l3.east) -- (zone_l3.west);

    % Flechas verticales de transición de paradigma (I4.0 -> Gemelo -> I5.0)
    \draw[flow spine] (sp_l1.north) -- node[edge lbl, align=center, fill=white, inner sep=1.5pt] {Despliegue\\Edge} (sp_l2.south);
    \draw[flow spine] (sp_l2.north) -- node[edge lbl, align=center, fill=white, inner sep=1.5pt] {Enfoque\\Humano} (sp_l3.south);

    % =========================================================================
    % 5. ENCABEZADO GLOBAL Y LEYENDA INFERIOR (CENTRADOS Y ALINEADOS AL ANCHO)
    % =========================================================================
    \coordinate (hdr_y) at ($(zone_l3.north) + (0, 0.54cm)$);
    \coordinate (hdr_l) at (sp_l3.west |- hdr_y);
    \coordinate (hdr_r) at (zone_l3.east |- hdr_y);

    \node[font=\sffamily\small\bfseries, text=archSlateDark, inner ysep=2pt]
        (hdr_txt) at ($(hdr_l)!0.5!(hdr_r)$) {%
        AgroMolinos del Valle S.A. \;---\; Arquitectura Integral de Transición Industria 4.0 $\rightarrow$ Industria 5.0%
    };

    \begin{pgfonlayer}{background}
        \node[
            rectangle, rounded corners=3.5pt,
            draw=archSlateLine, line width=0.8pt,
            fill=archPanelBg,
            inner xsep=0pt, inner ysep=4pt,
            fit=(hdr_l)(hdr_r)(hdr_txt)
        ] (hdr_global) {};
    \end{pgfonlayer}

    % Leyenda compacta inferior perfectamente distribuida y centrada al ancho total
    \coordinate (leg_y) at ($(zone_l1.south) + (0, -0.56cm)$);
    \coordinate (leg_l) at (sp_l1.west |- leg_y);
    \coordinate (leg_r) at (zone_l1.east |- leg_y);

    \node[edge lbl, anchor=west] (leg_t1) at ($(leg_l)!0.052!(leg_r)$) {Transferencia offline I4.0 (parámetros y pesos)};
    \draw[flow amber] ($(leg_t1.west) + (-0.68cm, 0)$) -- ($(leg_t1.west) + (-0.06cm, 0)$);

    \node[edge lbl, anchor=west] (leg_t2) at ($(leg_l)!0.432!(leg_r)$) {Bus reactivo 10\,Hz};
    \draw[flow blue bi] ($(leg_t2.west) + (-0.68cm, 0)$) -- ($(leg_t2.west) + (-0.06cm, 0)$);

    \node[edge lbl, anchor=west] (leg_t3) at ($(leg_l)!0.635!(leg_r)$) {Operación diegética I5.0 / SCADA};
    \draw[flow green bi] ($(leg_t3.west) + (-0.68cm, 0)$) -- ($(leg_t3.west) + (-0.06cm, 0)$);

    \node[edge lbl, anchor=west] (leg_t4) at ($(leg_l)!0.902!(leg_r)$) {Auditoría QA};
    \draw[flow green uni] ($(leg_t4.west) + (-0.68cm, 0)$) -- ($(leg_t4.west) + (-0.06cm, 0)$);

    \begin{pgfonlayer}{background}
        \node[
            rectangle, rounded corners=3pt,
            draw=archSlateLine, line width=0.75pt,
            fill=archPanelBg,
            inner xsep=0pt, inner ysep=4.5pt,
            fit=(leg_l)(leg_r)(leg_t1)(leg_t4)
        ] (legend_box) {};
    \end{pgfonlayer}
\end{tikzpicture}%
}
\caption{Arquitectura integral en tres capas de la transición de Industria~4.0 a Industria~5.0 en AgroMolinos del Valle S.A. La capa base consolida el modelado electromecánico y térmico de los Avances~1 y~2 en MATLAB/Simulink junto con el entrenamiento de clasificadores MLP y SVM~RBF; la capa central articula el núcleo en tiempo real del \textit{Serious Game} sobre Godot~4.7.2 (Vulkan~1.3 Forward+) mediante \texttt{EventBus.gd} y \texttt{SimulationManager.gd} a \SI{10}{\hertz}; y la capa superior materializa la operación centrada en el humano mediante interacción diegética~3D, tablet SCADA nativa \texttt{[TAB]} y certificación QA automatizada.}
\label{fig:arquitectura_global}
\end{figure*}


La Fig.~\ref{fig:arquitectura_global} ilustra la arquitectura ciberfísica en tres capas que articula la transición desde la ingeniería fuera de línea de Industria~4.0 hasta la operación centrada en el humano de Industria~5.0. En la capa central de tiempo real, el simulador se organiza en \num{16} escenas modulares \texttt{.tscn} (\texttt{main.tscn} más \num{15} subescenas distribuidas en entorno, equipos interactivos, infraestructura, jugador, cadena productiva, operarios e interfaz). El desacoplamiento entre física, IA, interacción y visualización se garantiza mediante dos \emph{autoloads} globales: \texttt{EventBus.gd}, que expone \num{16} señales troncales desacopladas (telemetría continua, inyección/resolución de fallas, alertas Edge AI, parada de emergencia, EPP, LOTO, zonas y notificaciones), y \texttt{Simulation\allowbreak Manager.gd}, responsable de coordinar el ciclo de telemetría a \qty{10}{\hertz}, el enclavamiento de arranque de maquinaria hasta completar la inducción, la parada de emergencia global y el bloqueo de reparación segura (\texttt{are\_\allowbreak machines\_\allowbreak stopped()}).

Para sostener fluidez interactiva sin sacrificar fidelidad visual, el motor se configuró en resolución nativa (\texttt{scaling\_3d/scale=1.0}) con \emph{antialiasing} MSAA~4x y FXAA, oclusión ambiental en espacio de pantalla (SSAO), resplandor (\emph{Glow}), búfer de profundidad \emph{Reverse-Z} de \num{32} bits y mapas de sombras direccionales PSSM de \num{4096}\,px. Sobre esta base actúa \texttt{Plant\allowbreak Visibility\allowbreak Optimizer.gd}, que elimina caras coplanares duplicadas, convierte materiales con \texttt{alphaMode=MASK} innecesario a opacos con \emph{Early-Z}, aplica \emph{culling} sectorial entre el bloque administrativo y la nave de producción, y agrupa la arquitectura estática en celdas espaciales LOD de \qty{16}{\meter}. En paralelo, \texttt{Plant\allowbreak Spatial\allowbreak Builder.gd} construye el entorno físico y normativo de toda la parcela.

El simulador aplica una \textbf{política estricta de cero primitivas 3D procedurales}: toda la geometría visible proviene de activos descargados con mallas originales de alta densidad y texturas en alta definición. En concreto, la infraestructura completa reside en \texttt{planta\_\allowbreak completa.glb} (\num{1338166} triángulos y \num{69} texturas originales de hasta $2048\times 1536$\,px), el Centro de Control de Motores en \texttt{panel\_\allowbreak control.glb} (\num{42140} triángulos y \num{20} texturas de $1024\times 1024$\,px) y la dotación humana en seis modelos \texttt{trabajador\_*.glb} (\num{7000} triángulos por operario con mapas difusos, ORM y de normales).

En el dominio de física espacial, \texttt{Plant\allowbreak Spatial\allowbreak Builder.gd} habilita el suelo industrial antideslizante \texttt{Suelo\allowbreak FisicoPBR} con coeficiente de fricción de \num{0.85} y genera exactamente \num{846} colisionadores convexos \texttt{Box\allowbreak Shape3D} sobre muros, losas, pasarelas, silos y maquinaria, con \textbf{cero} mallas cóncavas (\texttt{Concave\allowbreak Polygon\allowbreak Shape3D}), asegurando estabilidad numérica en CPU y tránsito fluido en primera persona. Asimismo, identifica y vuelve operables \num{28} puertas y portones interactivos (\texttt{Interactive\allowbreak Door.gd}): las \num{23} puertas batientes y dobles se alinean automáticamente al plano exacto del marco mediante Análisis de Componentes Principales bidimensional (PCA~2D, Ec.~\eqref{eq:pca_eigen}), logrando giros limpios de \qty{90}{\degree} sin penetración de jambas, complementadas por cortinas enrollables y portones corredizos automáticos. Sobre el piso se proyecta la demarcación normativa industrial sin artefactos de \emph{Z-fighting}: Zona~1 de Molienda en franjas amarillo/negro conforme a OSHA~1910.144~\cite{OSHA1910}; Zona~2 del tablero CCM \qty{460}{\volt} en perímetro rojo de despeje de \qty{1.0}{\meter} según RETIE Art.~20 y NFPA~70E~\cite{RETIE2026,NFPA70E}; pasillo peatonal seguro en verde; zona de recepción en azul; almacén de sacos en blanco; y Zona~ATEX~22 para polvo combustible de cereal en naranja según NFPA~61~\cite{NFPA61}.

\begin{figure*}[!t]
\centering
\begin{minipage}[t]{0.488\textwidth}
\centering
\includegraphics[width=\textwidth]{godot_planta_exterior.png}
\par\vspace{2pt}{\footnotesize (a) Vista aérea exterior del complejo industrial y circuitos vehiculares.}
\end{minipage}\hfill
\begin{minipage}[t]{0.488\textwidth}
\centering
\includegraphics[width=\textwidth]{godot_planta_interior_motor.png}
\par\vspace{2pt}{\footnotesize (b) Nave de molienda, tablero CCM \qty{460}{\volt}, NPCs y HUD minimalista.}
\end{minipage}
\caption{Entorno tridimensional de AgroMolinos del Valle S.A. en Godot~4.7.2 (Vulkan~1.3 Forward+): (a)~complejo completo con camión amarillo de descarga sobre la rampa de concreto del Patio Norte demarcada con perímetro OSHA, dos tractocamiones articulados en muelles de despacho y línea guía peatonal verde; (b)~interior de la nave frente al CCM \qty{460}{\volt} con zona roja RETIE, molino de martillos al fondo, Supervisor SST Carlos Mendoza aplicando bloqueo verbal ante ingreso sin EPP, Técnico Electromecánico NPC y elementos del HUD en primera persona.}
\label{fig:godot_planta_3d}
\end{figure*}

La dinámica de planta se cierra mediante \texttt{Animated\allowbreak Production\allowbreak Logistics.gd}, que sincroniza las \num{15} máquinas de la línea (rotación longitudinal del tornillo sinfín sobre el eje~$X$, flujos de partículas de grano y harina en tolva, trituradora y ciclones, y ensacado de \qty{50}{\kilogram}) directamente con las RPM instantáneas del motor de \qty{10}{HP}. En los circuitos exteriores pavimentados operan tractocamiones articulados con física de quinta rueda y giro natural acotado a $\le \qty{42}{\degree}$, cinemática exacta de dos ejes sobre la rampa de concreto del Patio Norte y apertura automática de los portones corredizos de parcela. El cargue manual de bultos de \qty{50}{\kilogram} por parte de los operarios de logística y muelle utiliza locomoción biomecánica dedicada (\texttt{Human\allowbreak Gait\allowbreak Animator.gd}, modo \texttt{carry\_walk}) y se ejecuta exclusivamente cuando el tractocamión permanece físicamente estacionado frente al muelle. Toda la sonorización 2D/3D se genera en memoria mediante flujos PCM normalizados (\texttt{Industrial\allowbreak Sound\allowbreak Synthesizer.gd}), abarcando zumbido electromagnético del motor, arco eléctrico de \qty{460}{\volt}, conmutación de \emph{breakers}, caudal de cereal, motores diésel con alarma de reversa y alertas de la Tablet. La Fig.~\ref{fig:godot_planta_3d} evidencia tanto la implantación exterior completa como la operación interna frente al CCM \qty{460}{\volt}.

\subsection{Flujo Pedagógico Industria 5.0: Inducción RETIE, LOTO y Rúbrica}\label{subsec:flujo_pedagogico}

En concordancia con la evidencia reciente sobre la eficacia del entrenamiento inmersivo en seguridad eléctrica~\cite{Stefan2023,Stefan2024}, el recorrido pedagógico (\texttt{Tutorial\allowbreak Induction\allowbreak Manager.gd}) impone una experiencia realista libre de menús invasivos. El estudiante inicia en la Portería exterior; ningún modal ni curso SST se despliega automáticamente al arrancar la simulación. La inducción comienza únicamente cuando el jugador se aproxima e interactúa (\texttt{E} o clic) con el \textbf{Supervisor SST (Carlos Mendoza)}, quien relata la historia, actividad productiva, misión y visión de AgroMolinos del Valle S.A. antes de ofrecer la capacitación normativa e iniciar el recorrido guiado.

El primer filtro obligatorio consiste en ingresar a los Vestidores (\texttt{Locker\allowbreak EPP\allowbreak Station.gd}) para dotarse de \textbf{6/6 prendas EPP reglamentarias}~\cite{RETIE2026,NFPA70E,OSHA1910}: (1)~casco dieléctrico Clase~E (\qty{20}{\kilo\volt}), (2)~gafas de seguridad ANSI~Z87.1, (3)~protección auditiva (NRR \qty{27}{\decibel}), (4)~guantes dieléctricos Clase~0 (\qty{1000}{\volt}), (5)~botas dieléctricas ASTM~F2413 (\qty{18}{\kilo\volt}) y (6)~ropa ignífuga NFPA~70E Categoría~2 ($8\,\mathrm{cal/cm^2}$). Mientras falte dotación, el área de enclavamiento de acceso detecta la intrusión a la nave industrial, descuenta la sanción correspondiente, actualiza el avatar 3D del HUD (\texttt{SIN EPP} $\leftrightarrow$ \texttt{EPP 6/6 OK}) y desplaza al Supervisor SST para interceptar al operario y ordenarle regresar a vestidores, tal como se observa en la Fig.~\ref{fig:godot_planta_3d}(b).

Antes de intervenir cualquier componente energizado en el tablero CCM \qty{460}{\volt}, el módulo \texttt{Loto\allowbreak Golden\allowbreak Rules\allowbreak Controller.gd} exige ejecutar en secuencia estricta $1\to 2\to 3\to 4\to 5$ las \textbf{5 Reglas de Oro del RETIE}~\cite{RETIE2026,Res5018}: (1)~corte visible del interruptor principal, (2)~bloqueo y tarjeteo con candado LOTO, (3)~verificación de ausencia de tensión (\qty{0}{\volt}) mediante multímetro, (4)~instalación de puesta a tierra temporal en cortocircuito y (5)~señalización perimetral; cualquier intento fuera de orden es rechazado inmediatamente con aviso preventivo.

El recorrido guiado comprende \textbf{20 estaciones verificadas}, desde la Estación~1 (Portería) y la Estación~2 (Vestidores EPP), atravesando las zonas de recepción, tolva, trituradora de \qty{10}{HP}, CCM \qty{460}{\volt}, secador, ciclones ATEX, almacén y patio térmico exterior, hasta la Estación~20 de cierre y entrega de consignas operativas frente al tablero CCM \qty{460}{\volt}. La trayectoria peatonal se dibuja sobre el suelo mediante \texttt{Tour\allowbreak Floor\allowbreak Guide\allowbreak Line} a lo largo de \num{61} \emph{waypoints} trazados sobre aceras, rampas de concreto y pasillo peatonal verde, garantizando cero cruces por muros o zonas de peligro. Durante el turno interactúan \num{6} operarios NPC autónomos (\texttt{WorkerNPC.gd}): Supervisor SST (Carlos Mendoza), Técnico Electromecánico (Andrés Rojas), Operario Logístico (Luis Herrera), Jefe de Planta (Ing.~Diana Morales, quien entrega reporte vivo del estado y fallas de planta) y dos señaleros/operarios de muelle (Jorge Gómez y Mateo Silva).

El desempeño seguro se cuantifica en tiempo real mediante \texttt{Safety\allowbreak Evaluation\allowbreak Scorecard.gd} en una escala académica continua de \num{0.0} a \num{5.0} (nota base \num{5.0}). Las infracciones verificadas aplican deducciones inmediatas: \num{-0.8} por acceso sin EPP, \num{-1.0} por tránsito sin botas dieléctricas, \num{-0.6} por ingreso a molienda sin protección auditiva, \num{-1.5} por invasión del perímetro rojo del CCM sin guantes, \num{-2.5} por intentar intervenir el molino en giro y \num{-3.0} por cortar un conductor energizado sin LOTO. En contrapartida, el cumplimiento otorga bonificaciones verificadas: \num{+1.0} por completar en secuencia las 5 Reglas de Oro LOTO, \num{+1.0} por diagnosticar y remediar físicamente una falla y \num{+0.5} por certificar la inducción completa. Al finalizar la práctica, la bitácora completa y la nota definitiva se exportan en formato JSON estructurado (\texttt{Reporte\_\allowbreak Estudiante\_\allowbreak RETIE.json}) como evidencia evaluativa auditable.

\subsection{Gemelo Físico en Tiempo Real, Edge AI y 5 Fallas Diegéticas}\label{subsec:gemelo_ia_fallas}

El acoplamiento multifísico e inteligente corre nativamente en GDScript~2.0: el lazo físico \texttt{\_physics\_process} integra las ecuaciones diferenciales del motor a \qty{50}{\hertz}, mientras la emisión del vector sensorial completo hacia la IA y la Tablet ocurre a \qty{10}{\hertz} mediante \texttt{EventBus.\allowbreak telemetry\_\allowbreak vector\_\allowbreak updated}.

\paragraph*{Gemelo Electromecánico, Térmico y Vibratorio}
La clase \texttt{Induction\allowbreak Motor\allowbreak Kloss.gd} implementa en tiempo real el motor WEG W22 IE3 de \qty{10}{HP} (\qty{7.457}{\kilo\watt}, \qty{460}{\volt}, \qty{60}{\hertz}, \num{4} polos, $I_{\text{FLA}}=\qty{13.5}{\ampere}$, $I_{\text{LRA}}=\qty{81.0}{\ampere}$, $T_{\max}=\qty{101.75}{\newton\meter}$, $s_{\max}=\num{0.15}$)~\cite{WEG2024,IEC60034}. El par electromagnético sigue la ecuación de Kloss escalada por relación $V/f$ (Ec.~\eqref{eq:kloss_vf}), incorporando compensación óhmica $IR$ (\emph{voltage boost}) de \qty{18}{\volt} en baja frecuencia ($f<\qty{15}{\hertz}$) para preservar el par de despegue:
\begin{equation}
\begin{split}
V_{\text{ap}} &= \min\!\Big(V_{LL},\; V_{LL}\min\!\big(1, f/f_n\big) \\
&\qquad\qquad + 18\operatorname{clip}_{[0,1]}\!\big(1 - f/15\big)\Big).
\end{split}
\label{eq:godot_vf_ir}
\end{equation}
La velocidad rotórica evoluciona por la ecuación de Newton--Euler (Ec.~\eqref{eq:newton_euler}) con inercia combinada $J_{\text{tot}}=\qty{0.35}{\kilogram\meter\squared}$ (\qty{0.045}{\kilogram\meter\squared} del motor más \qty{0.305}{\kilogram\meter\squared} del molino de martillos) y fricción viscosa $B=\qty{0.012}{\newton\meter\second\per\radian}$, oponiéndose al par resistente dependiente del flujo másico de cereal $\dot{m}$ más un sobrecargo de \qty{120}{\newton\meter} ante atascamiento. En el dominio térmico, la red de Cauer de primer orden (Ec.~\eqref{eq:red_cauer}) integra las pérdidas Joule estatóricas $(I_a^2+I_b^2+I_c^2)R_s$, pérdidas rotóricas y pérdidas en el hierro (\qty{180}{\watt}) sobre una capacitancia térmica $C_{\text{th}}=\qty{185}{\joule\per\kelvin}$, con resistencia térmica convectiva $R_{\text{th}}=\qty{0.068}{\kelvin\per\watt}$ en ventilación sana (modulada por la velocidad del ventilador acoplado al eje) frente a $R_{\text{th}}=\qty{0.210}{\kelvin\per\watt}$ bajo rejilla obstruida~\cite{Thermal2025,Chapman2012}. Complementariamente, el módulo de vibración evalúa $v_{\text{rms}}$ y $v_{\text{peak}}$ conforme a ISO~10816-3/20816-3~\cite{ISO10816,ISO20816}, partiendo de una línea base en Zona~A ($\approx\qty{1.10}{\milli\meter\per\second}$) que escala a Zona~D ($>\qty{4.5}{\milli\meter\per\second}$) por pulsación magnética a \qty{120}{\hertz} en pérdida de fase o por excentricidad mecánica ($>\qty{5.5}{\milli\meter\per\second}$) ante pernos flojos.

\paragraph*{Optimización Energética VFD frente a DOL}
El controlador \texttt{VFDDrive\allowbreak Controller.gd} permite conmutar en vivo entre accionamiento por variador (\textbf{VFD}) y arranque directo (\textbf{DOL}). En modo VFD aplica una rampa suave de \qty{6}{\hertz\per\second} (\qty{10}{\second} de $0$ a \qty{60}{\hertz}) que limita la corriente de arranque a $\le 1.3\,I_{\text{FLA}}$ ($\approx\qty{17.5}{\ampere}$, frente a \qty{81.0}{\ampere} en DOL) y ejecuta un lazo PID de demanda proporcional al flujo másico $\dot{m}$: regula a \qty{42}{\hertz} para carga parcial ($\dot{m}\le\qty{0.35}{\kilogram\per\second}$), interpola linealmente en el rango intermedio y alcanza \qty{60}{\hertz} a plena carga ($\dot{m}\ge\qty{0.80}{\kilogram\per\second}$). Aplicando la ley cúbica de afinidad y el factor de emisión colombiano de la UPME (\qty{0.164}{\kilogram\per\kilo\watt\hour} de $\text{CO}_2\text{e}$, Ec.~\eqref{eq:afinidad_upme})~\cite{UPME2024,Resa2019}, cuantifica en vivo el consumo específico (\unit{\kilo\watt\hour\per\tonne}) y las emisiones evitadas frente a la línea base DOL (\qty{7.52}{\kilo\watt\hour\per\tonne} frente a \qty{10.50}{\kilo\watt\hour\per\tonne}).

\paragraph*{Inferencia Edge AI 100\% Nativa}
El motor \texttt{EdgeAIEngine.gd} ejecuta inferencia completamente local en GDScript a \qty{10}{\hertz}, sin depender de MATLAB ni servidores externos~\cite{EdgeCloud2022,AI2026}. Incorpora literalmente los \textbf{316 parámetros entrenados} ($W_1\in\mathbb{R}^{8\times 16}$, $\mathbf{b}_1\in\mathbb{R}^{16}$, $W_2\in\mathbb{R}^{16\times 8}$, $\mathbf{b}_2\in\mathbb{R}^{8}$, $W_3\in\mathbb{R}^{8\times 4}$, $\mathbf{b}_3\in\mathbb{R}^{4}$) y los vectores de estandarización $\boldsymbol{\mu},\boldsymbol{\sigma}$ extraídos de \texttt{modelo\_\allowbreak ia\_\allowbreak entrenado.mat} del Avance~2. Para conciliar la corriente nominal de placa del motor IE3 (\qty{13.5}{\ampere}) con el dominio p.u.\ de entrenamiento (\qty{9.363}{\ampere}), el preprocesador aplica una adaptación de dominio lineal antes de la normalización $Z$-Score (Ec.~\eqref{eq:zscore}):
\begin{equation}
I_{\text{MLP}} = I_{\text{rms}} \cdot \frac{9.363}{13.5}.
\label{eq:godot_adaptacion}
\end{equation}
Tras la propagación hacia adelante de la red MLP $8\to 16\to 8\to 4$ (ReLU y Softmax, Ec.~\eqref{eq:mlp_softmax}), se realiza validación cruzada paralela mediante SVM con kernel RBF ($\gamma=\num{0.125}$, $C=\num{10.0}$, Ec.~\eqref{eq:kernel_rbf})~\cite{Gangsar2020,Soualhi2022} reforzada con salvaguardas normativas RETIE/ISO para transitorios físicos (p.ej., $\text{IUF}\ge\qty{15}{\percent}$, $V_{\text{rms}}<\qty{414}{\volt}$ o $v_{\text{rms}}\ge\qty{3.2}{\milli\meter\per\second}$). Finalmente calcula en vivo el índice de salud MHI (Ec.~\eqref{eq:mhi}, acotado a [\qty{5.0}{\percent}, \qty{99.6}{\percent}]) y la vida útil remanente RUL (Ec.~\eqref{eq:rul}, combinando envejecimiento térmico de Arrhenius y fatiga SKF $L_{10h}$ hasta \qty{24000}{\hour}), emitiendo alertas push con causa raíz, pronóstico y lista de verificación RCM~\cite{Moubray1997}.

\begin{table*}[!t]
\begingroup\setlength{\tabcolsep}{3.5pt}\renewcommand{\arraystretch}{1.08}
\centering\footnotesize
\caption{Matriz de las cinco fallas físicas diegéticas implementadas en el mundo 3D (cero botones 2D de inyección de falla).}
\label{tab:fallas_diegeticas}
\begin{tabularx}{\textwidth}{@{} c >{\raggedright\arraybackslash}p{2.3cm} >{\raggedright\arraybackslash}p{2.65cm} >{\raggedright\arraybackslash}p{3.1cm} >{\raggedright\arraybackslash}X >{\raggedright\arraybackslash}p{3.25cm} @{}}
\toprule
\textbf{\#} & \textbf{Falla diegética} & \textbf{Script 3D} & \textbf{Herramienta e inducción} & \textbf{Respuesta electromecánica e IA} & \textbf{Reparación segura (parada previa)} \\
\midrule
1 & Corte físico fase~L3 & \texttt{Physical\allowbreak Cable\allowbreak Conductor.gd} & \texttt{2}~Alicates: cortar conductor \texttt{L3\_\allowbreak Azul} (si no hay LOTO: arco \qty{460}{\volt} y sanción \num{-3.0}). & $I_c=\qty{0}{\ampere}$, $I_a,I_b\to\sqrt{3}\,I_{\text{carga}}$, $\text{IUF}>\qty{25}{\percent}$, $\text{THD}\approx\qty{13.5}{\percent}$, $v_{\text{rms}}>\qty{4.5}{\milli\meter\per\second}$ $\to$ \texttt{FALLA\_\allowbreak ELECTRICA}. & Detener máquinas (Tablet/LOTO) y empalmar fase~L3 con \texttt{5}~Kit~LOTO. \\
2 & Atasco en tolva por sobrealimentación & \texttt{Physical\allowbreak Hopper\allowbreak Feeder.gd} + \texttt{Interactive\allowbreak GrainBag.gd} & \texttt{1}~Manos: arrojar $\ge 2$ sacos de \qty{50}{\kilogram}; o \texttt{3}~Llave~3/4'': forzar compuerta al \qty{160}{\percent}. & $\dot{m}=\qty{4.5}{\kilogram\per\second}$, $T_L>T_{\max}$, rotor bloqueado ($N_r<\num{50}$\,RPM), $I_{\text{rms}}\to I_{\text{LRA}}$, polvo ATEX $\to$ \texttt{SOBRE\allowbreak CALENTAMIENTO}. & Parada obligatoria en Tablet y desatasco con \texttt{3}~Llave~3/4'' (flujo a \qty{0.80}{\kilogram\per\second}). \\
3 & Bajón de tensión a \qty{340}{\volt} en CCM & \texttt{Physical\allowbreak Breaker\allowbreak Panel.gd} & Abrir gabinete \texttt{panel\_\allowbreak control.glb} y bajar palanca \texttt{Breaker\_\allowbreak L3} con \texttt{2}~Alicates o \texttt{3}~Llave~3/4''. & $V_{LL}:\qty{460}{\volt}\to\qty{340}{\volt}$, caída cuadrática de par ($T_e\propto V^2$), sobrecorriente compensatoria y $\text{THD}\approx\qty{9.2}{\percent}$ $\to$ \texttt{FALLA\_\allowbreak ELECTRICA}. & Parada previa en Tablet y rearme de palanca a \qty{460}{\volt} con \texttt{2}~Alicates o \texttt{3}~Llave~3/4''. \\
4 & Obstrucción de rejilla de ventilación & \texttt{Ventilation\allowbreak Cover.gd} & \texttt{2}~Alicates aislados: colocar obstrucción física sobre la rejilla del ventilador trasero. & $R_{\text{th}}:\qty{0.068}{\kelvin\per\watt}\to\qty{0.210}{\kelvin\per\watt}$, gradiente $dT_w/dt$ acelerado y deriva $T_w>\qty{105}{\celsius}$ $\to$ \texttt{SOBRE\allowbreak CALENTAMIENTO}. & Parada obligatoria en Tablet y retiro del obstáculo con \texttt{2}~Alicates. \\
5 & Soltura de pernos de bancada & \texttt{Mechanical\allowbreak Bolt.gd} & \texttt{3}~Llave~3/4'': aflojar perno de anclaje (\texttt{Bolt\_\allowbreak Front\_\allowbreak Right}) en bancada del motor. & Excentricidad estructural 1X/2X, factor de cresta \num{2.5}, $v_{\text{rms}}>\qty{5.5}{\milli\meter\per\second}$ (Zona~D ISO) $\to$ \texttt{FALLA\_\allowbreak MECANICA}. & Parada obligatoria en Tablet y reapriete de torque con \texttt{3}~Llave~3/4''. \\
\bottomrule
\end{tabularx}
\endgroup
\end{table*}

\paragraph*{Regla Cardinal de Fallas Diegéticas}
El simulador suprime por completo cualquier botón 2D de inyección de fallas en la interfaz. Toda perturbación y remediación se ejecuta corporizadamente en el espacio 3D mediante el cinturón de cinco herramientas (\texttt{1}~Manos, \texttt{2}~Alicates dieléctricos \qty{1000}{\volt}, \texttt{3}~Llave fija 3/4'', \texttt{4}~Multímetro CAT~III y \texttt{5}~Kit~LOTO). Además, cada rutina de reparación consulta obligatoriamente el enclavamiento \texttt{Simulation\allowbreak Manager.\allowbreak are\_\allowbreak machines\_\allowbreak stopped()}: si el motor permanece energizado o en giro, la intervención es bloqueada y el sistema exige detener primero la planta desde la Tablet (\texttt{[TAB]}) o aislar con LOTO. La Tabla~\ref{tab:fallas_diegeticas} consolida el comportamiento técnico de las cinco fallas diegéticas.

\subsection{Interfaz In-Game: HUD Minimalista y Tablet Industrial SCADA}\label{subsec:interfaz_ingame}

Para preservar la inmersión sin saturar el campo visual, la interfaz de usuario se divide en un HUD minimalista de primera persona y una Tablet Industrial SCADA desplegable bajo demanda con \texttt{[TAB]}.

\paragraph*{HUD de Primera Persona}
Gestionado por \texttt{HUD\allowbreak Controller.gd} y \texttt{HUDHotbar\allowbreak AndAvatar.gd}, el HUD visible en la Fig.~\ref{fig:godot_planta_3d}(b) incorpora seis elementos estrictamente funcionales y sin solapamientos: (1)~una mira central fija blanca (\texttt{+}, $14\times 2$\,px) con borde negro de \qty{1}{px} que mantiene contraste constante sobre cualquier superficie; (2)~una píldora superior que anuncia la zona actual del operario (p.ej., \texttt{Pasillo Peatonal}); (3)~la barra inferior de cinco herramientas (\texttt{1}--\texttt{5}) con resaltado activo en dorado y cambio por teclado o rueda del ratón; (4)~un avatar 3D reactivo a la izquierda de la barra que refleja en vivo la vestimenta y la insignia \texttt{SIN EPP} $\leftrightarrow$ \texttt{EPP 6/6 OK}; (5)~una caja de subtítulos naturales para los diálogos de NPCs (\texttt{Nombre (Cargo): ...}); y (6)~una columna inferior derecha unificada que apila las notificaciones críticas de seguridad y avisos de tablet directamente encima de la tarjeta compacta de controles (\texttt{WASD}, \texttt{Shift}, \texttt{E/Clic}, \texttt{1-5}, \texttt{TAB}, \texttt{F11}).

\paragraph*{Tablet Industrial SCADA Nativa (\texttt{[TAB]})}
Implementada en \texttt{Industrial\allowbreak Tablet\allowbreak Controller.gd} y siete módulos auxiliares (\texttt{Tablet*.gd}), la Tablet es \textbf{100\% nativa dentro de Godot} sin depender de navegadores web ni servidores locales. Conserva un tamaño geométrico invariante de $1228\times 748$\,px en chasis y $1180\times 700$\,px en pantalla útil, con apertura y cierre directos mediante \texttt{TAB} o \texttt{ESC}. Su encabezado fijo presenta el título \texttt{AGROMOLINOS SCADA}, una píldora de estado operacional en vivo (\texttt{Normal}, \texttt{Alerta}, \texttt{Falla crítica} o \texttt{Parada}), el indicador continuo de calificación académica \texttt{RETIE} y el pulsador rojo exclusivo \texttt{Parar Máquinas} / \texttt{Reanudar Máquinas}, única parada de emergencia global del simulador. Su navegación interna se estructura en \textbf{nueve pestañas especializadas}, sintetizadas en la Tabla~\ref{tab:tablet_pestanas}.

\begin{table*}[!t]
\begingroup\setlength{\tabcolsep}{3.6pt}\renewcommand{\arraystretch}{1.06}
\centering\footnotesize
\caption{Estructura funcional de las nueve pestañas de la Tablet Industrial SCADA nativa (\texttt{[TAB]}) en Godot~4.7.2.}
\label{tab:tablet_pestanas}
\begin{tabularx}{\textwidth}{@{} c >{\raggedright\arraybackslash}p{2.35cm} >{\raggedright\arraybackslash}p{4.35cm} >{\raggedright\arraybackslash}X @{}}
\toprule
\textbf{\#} & \textbf{Pestaña} & \textbf{Módulos / controles clave} & \textbf{Alcance operativo, analítico y normativo in-game} \\
\midrule
1 & \texttt{Monitoreo IoT} & Estado CCM/VFD, \texttt{Modo: VFD/DOL}, audio, banner de alertas, filtros de zona, \texttt{Tablet\allowbreak Oscilloscope\allowbreak Chart.gd} y \emph{Blackbox}. & Supervisa en vivo los \num{16} activos (\num{15} máquinas en orden de proceso + CCM \qty{460}{\volt}), con oscilogramas nativos de corrientes $I_a/I_b/I_c$, temperatura $T_w$, vibración ISO y caja negra cronológica. \\
2 & \texttt{Cámaras} & \num{24} visores 3D \texttt{SubViewport} en \num{6} grupos (\texttt{Recepción}, \texttt{Molienda}, \texttt{Ciclones}, \texttt{Aire y CCM}, \texttt{Patio}, \texttt{Operarios}) con \texttt{Iso}, \texttt{Frente}, \texttt{Top}, \texttt{Rotar}. & Inspección visual remota en vivo con telemetría superpuesta y renderizado bajo demanda; se complementa en planta con el banco de \num{4} monitores CCTV en el Mezzanine (\texttt{Mezzanine\allowbreak CCTV\allowbreak Controller.gd}). \\
3 & \texttt{RETIE} & Tarjetas de 5 Reglas de Oro LOTO, fichas Motor \qty{10}{HP} / CCM-VFD / EPP (\texttt{0/6}--\texttt{6/6}) y tabla de riesgos. & Seguimiento en vivo de la consignación LOTO, verificación de dotación EPP y consulta rápida de los cinco riesgos críticos RETIE/ATEX con sus controles de ingeniería. \\
4 & \texttt{Operarios} & Ficha del jugador, nómina en turno y formulario de registro (nombres, edad, RH, cargo). & Registra personal autorizado en turno y enlaza la estación actual, el estado EPP y la nota de seguridad del estudiante. \\
5 & \texttt{Riesgos} & Evaluador interactivo por sistema y categoría de peligro (\texttt{Tablet\allowbreak Secondary\allowbreak Tabs\allowbreak Builder.gd}). & Analiza peligros eléctricos, mecánicos, térmicos y de polvo combustible ATEX Zona~22 con causas, consecuencias y salvaguardas. \\
6 & \texttt{Fisiología IEC} & Simulador log--log IEC~60479-1~\cite{IEC60479} (\texttt{Tablet\allowbreak Analytical\allowbreak Canvases.gd}) y cálculo energético. & Evalúa en tiempo real las curvas frontera AC-1 a AC-4 (Ec.~\eqref{eq:iec60479_interp}) y la energía corporal $E=I^2 R_b t$ (Ec.~\eqref{eq:energia_choque}) ante choques de corriente y duración configurables. \\
7 & \texttt{Hardware \& IA} & Tarjetas Edge AI (\num{8} vars, MLP, SVM RBF, SCADA), diagnóstico vivo, barra MHI, RUL, checklist RCM y diagrama de bloques. & Muestra certezas MLP/SVM, causa raíz, pronóstico, prescripción RCM interactiva y separación galvánica entre Etapa Control/Edge AI (\qty{3.3}{\volt}/\qty{5}{\volt}~DC) y Potencia Trifásica (\qty{460}{\volt}~AC). \\
8 & \texttt{Rúbrica} & Calificación RETIE (\texttt{0.0--5.0}), cuadro de ponderación, historial de eventos y grilla \texttt{E1}--\texttt{E19}. & Audita en vivo penalizaciones, bonificaciones y avance paso a paso por las estaciones del recorrido seguro (con cierre en CCM). \\
9 & \texttt{Historial} & Filtros por severidad (\texttt{Todas}, \texttt{Críticas}, \texttt{Alertas}, \texttt{Éxito}, \texttt{Info}), \num{4} KPIs contadores y lista persistente. & Conserva cronológicamente todas las notificaciones emitidas durante la sesión con limpieza controlada y conteo por criticidad. \\
\bottomrule
\end{tabularx}
\endgroup
\end{table*}

\begin{figure*}[!t]
\centering
\begin{minipage}[t]{0.485\textwidth}
\centering
\includegraphics[width=0.94\textwidth]{godot_tablet_iot.png}
\par\vspace{2pt}{\footnotesize (a) Pestaña \texttt{Monitoreo IoT}: alerta RETIE activa y enclavamiento en \qty{0}{\hertz}.}
\end{minipage}\hfill
\begin{minipage}[t]{0.485\textwidth}
\centering
\includegraphics[width=0.94\textwidth]{godot_tablet_camaras.png}
\par\vspace{2pt}{\footnotesize (b) Pestaña \texttt{Cámaras}: mosaico $2\times 2$ de visores 3D en zona \texttt{Recepción}.}
\end{minipage}
\caption{Capturas reales de las dos primeras pestañas de la Tablet Industrial SCADA (\texttt{[TAB]}): (a)~\texttt{Monitoreo IoT} desplegando la barra de nueve pestañas, el conmutador \texttt{Modo: VFD}, el botón \texttt{Audio: ON}, la alerta crítica roja \texttt{Acceso sin EPP prohibido (RETIE -0.8)} con enlace \texttt{Ver RCM} y el enclavamiento de seguridad con la línea detenida a \qty{0}{\hertz} y \qty{0.0}{\ampere} antes de completar la inducción; (b)~\texttt{Cámaras} renderizando simultáneamente cuatro visores \texttt{SubViewport} 3D (\texttt{01~Banda de Descarga}, \texttt{02~Banda Inclinada}, \texttt{03~Tolva Receptora} y \texttt{04~Banda Trituradora}) con telemetría superpuesta y controles orbitales \texttt{Iso}, \texttt{Frente}, \texttt{Top} y \texttt{Rotar}.}
\label{fig:godot_tablet_iot_cam}
\end{figure*}

\begin{figure*}[!t]
\centering
\begin{minipage}[t]{0.485\textwidth}
\centering
\includegraphics[width=0.94\textwidth]{godot_tablet_retie.png}
\par\vspace{2pt}{\footnotesize (a) Pestaña \texttt{RETIE}: 5 Reglas de Oro, estado \texttt{EPP: 0/6} y riesgos críticos.}
\end{minipage}\hfill
\begin{minipage}[t]{0.485\textwidth}
\centering
\includegraphics[width=0.94\textwidth]{godot_tablet_ia_falla.png}
\par\vspace{2pt}{\footnotesize (b) Pestaña \texttt{Hardware \& IA}: régimen base de espera segura (\texttt{MHI 99.6\%}).}
\end{minipage}
\caption{Capturas reales de supervisión normativa e inteligencia artificial embebida en la Tablet SCADA: (a)~pestaña \texttt{RETIE} con las cinco tarjetas secuenciales de consignación LOTO, fichas técnicas del \texttt{Motor 10 HP}, \texttt{Tablero CCM y VFD} y \texttt{Protección EPP} (\texttt{0/6}), y matriz de cinco riesgos eléctricos y ATEX; (b)~pestaña \texttt{Hardware \& IA} en estado base de espera segura (\texttt{0.0 A}, \texttt{IUF 0.0\%}, \texttt{460 V}, \qty{28.0}{\celsius}), mostrando la arquitectura nativa a \qty{10}{\hertz}, certezas \texttt{MLP 100.0\%} y \texttt{SVM 94.2\%}, \texttt{MHI 99.6\%}, \texttt{RUL 24000 h}, checklist RCM de tres pasos y esquema dual Control/Edge AI frente a Potencia Trifásica \qty{460}{\volt}.}
\label{fig:godot_tablet_retie_ia}
\end{figure*}

Las Figs.~\ref{fig:godot_tablet_iot_cam} a~\ref{fig:godot_tablet_rubrica} documentan el comportamiento real verificado de la interfaz durante una sesión operativa. En la Fig.~\ref{fig:godot_tablet_iot_cam}(a), la pestaña \texttt{Monitoreo IoT} exhibe el encabezado global con las nueve pestañas activas, el selector \texttt{Modo: VFD} y la alerta roja persistente generada tras el ingreso sin EPP (\texttt{RETIE -0.8}), mientras las tarjetas \texttt{01 Banda de Descarga} y \texttt{02 Banda Inclinada} permanecen enclavadas a \qty{0}{\percent} de velocidad y \qty{0.0}{\ampere} (\texttt{VFD 0 Hz}) hasta que el estudiante concluya su inducción. La Fig.~\ref{fig:godot_tablet_iot_cam}(b) muestra la pestaña \texttt{Cámaras} filtrada en el grupo \texttt{Recepción}, donde cuatro visores tridimensionales simultáneos proyectan en vivo las bandas de alimentación y la tolva receptora con lecturas de caudal y controles interactivos de ángulo de cámara.

Por su parte, la Fig.~\ref{fig:godot_tablet_retie_ia}(a) evidencia la pestaña \texttt{RETIE}, destacando el estado inicial limpio de las 5 Reglas de Oro, el contador \texttt{EPP: 0 / 6} en ámbar y la matriz de riesgos eléctricos críticos con sus barreras de control. En la Fig.~\ref{fig:godot_tablet_retie_ia}(b), la pestaña \texttt{Hardware \& IA} documenta el motor de inferencia nativo a \qty{10}{\hertz} operando en \textbf{estado base de espera segura} (\qty{0.0}{\ampere}, $\text{IUF}=\qty{0.0}{\percent}$, \qty{460}{\volt} y \qty{28.0}{\celsius}), donde la red MLP $8\text{--}16\text{--}8\text{--}4$ (\qty{100.0}{\percent}) y el clasificador SVM RBF (\qty{94.2}{\percent}) validan la condición saludable con $\text{MHI}=\qty{99.6}{\percent}$, $\text{RUL}=\qty{24000}{\hour}$ y las tres acciones preventivas del checklist RCM, sobre el diagrama de bloques que separa la etapa de control (\qty{3.3}{\volt}/\qty{5}{\volt}~DC) de la etapa de potencia trifásica (\qty{460}{\volt}~AC del molino \texttt{M-101}); la respuesta crítica de estos mismos indicadores bajo la inducción activa de las cinco fallas diegéticas se analiza y cuantifica rigurosamente en la Sección~\ref{sec:resultados}. Finalmente, la Fig.~\ref{fig:godot_tablet_rubrica} presenta la pestaña \texttt{Rúbrica}, donde la calificación en vivo refleja la deducción \texttt{4.2 / 5.0} registrada en el historial de infracciones (\num{-0.8} por ingreso sin EPP), acompañada de las reglas de ponderación y el tablero de avance por estaciones (\texttt{E1~OK}, \texttt{E2~Actual} hasta \texttt{E19}, antes del cierre operativo en la Estación~20 frente al CCM).

\begin{figure*}[!t]
\centering
\includegraphics[width=0.68\textwidth]{godot_tablet_rubrica.png}
\caption{Captura real de la pestaña \texttt{Rúbrica} en la Tablet SCADA: se observa la actualización instantánea de la nota normativa (\texttt{RETIE 4.2/5.0}) tras la deducción de \num{-0.8} por ingreso a la planta sin equipo de protección personal, el recuadro de ponderación de bonos y sanciones, y el tablero de progreso de las estaciones (\texttt{E1~OK}, \texttt{E2~Actual} hasta \texttt{E19}).}
\label{fig:godot_tablet_rubrica}
\end{figure*}
